Commentaires étoilés de MF :

%1. C’est peut-être parce que je ne suis pas du domaine mais je trouve que la facon de présenter ne clarifie rien, au contraire. Bref, il me semble que ce n’est pas autocontenu car on semble prendre pour acquis la connaissance du domaine pour que ce chapitre soit lisible. C’est contraire à l’esprit d’un mémoire ou d’une thèse. Personnellement, je ferais l’introduction du chapitre et je clarifierais après la démarche de présentation. Et je mettrais plus de détails sur celle-ci pour que ca soit clair pour quelqu’un hors domaine. Un genre de guide de lecture explicite. (fait)
% Maryse : À discuter avec Luc

%2. Sur l'identification du paragraphe : ca coupe la lecture, on sait pas a quoi ca sert. (fait)
%Commentaire de MF : Il faut réviser
% Maryse : oui

% Il faut télécharger le nouveau gabarit qui présente mieux les titres long (page titre) (fait)

%1. Structure

%a. Le chap. 1 devrait être mieux introduit, surtout l'aspect particulier de la présentation. Voir document annoté. (fait)

2. Langue

a. L’utilisation non justifiée de l’anglais dans certains cas : vérifier la vitrine linguistique (https://vitrinelinguistique.oqlf.gouv.qc.ca/) L'usage de l'italique ne suffit pas. Il faut utiliser la terminologie francophone, quitte à en suggérer au besoin. (trad. libre de ...) Exemples : figures et tableaux en annexe
% Maryse : 
b. Éliminer les précisions entre parenthèse de l'anglais triviales ou mot à mot
% Maryse : 
% Luc : Mettre tout en francais suivi d'une virgule et après en anglais.
% Maxime : Si je fais cela, le texte suivant ne sera pas en italique puisque déjà traduit demandé à MF.
c. Enlever les italiques un peu partout, normalement les italiques sont pour les termes en langue étrangère qui sont non traduits. Voir guide de rédaction.
% Luc : Regarder l'article de la vitrine sur le comportement à adopter sur la traduction.
% Maryse : https://vitrinelinguistique.oqlf.gouv.qc.ca/23207/la-redaction-et-la-communication/bibliographie-et-citations/citations/citation-en-langue-etrangere

% Maxime : Inversé les listes entre anglais et francais, demander précision à MF.
d. Les guillemets devraient être utilisés uniquement pour des citations dans le texte. Utiliser une police de caractère différente pour le code ou les mots-clés.
e. Les tables (des matières, figures, tableaux, ...) devraient être en français (Latex permet d'avoir un titre différent, voir gabarit).
% Luc : Faire un encadré avec le titre en anglais à l'intérieur et puis ajouté un titre en francais à la figure. S'il y a nécessité d'avoir plus d'information accompagner le titre de la figure avec un note de bas de page avec la traduction de l'information.
f. Traduire les titres et descriptifs de figures tirées de publications en anglais.
% Luc : Lorsque tu contenues de tableaux nécessites des tradutions conserver un tableau représentant la traduction dans l'annexe. Et ajouter une note de bas de page expliquant que la tracabilité de la traduction se trouve dans l'annexe.

3. Bibliographie : voir guide

%a. Tous les « récupéré de » devraient être retirés. C’est pertinent pour des site web, et non des publications scientifiques. Ne conserver que les DOI. (fait)
b. Certaines notices sont incomplètes.
% c. Les articles non publiés (ex. sur Xarchiv) doivent être mentionnés : non publié. (fait)
d. Les références multiples devraient être en ordre numériques
%e. Ne conserver que les DOI et enlever les URL redondants. (Fait)
%f. Ajouter les DOI avec hyperliens pour toute référence pertinente
g. Les mots principaux des titres anglais devraient débuter par des majuscules (mots principaux)
% Max : faire de son mieux, Question MF : 
% Quelle règle je suis parce que je dois avoir écrit sur les sites web. J'ai reproduit ce qui est adns l'exemple....

4. Éléments flottants (figures, tableaux, codes, algorithmes) numérotés

a. Listes (figures, tableaux, etc.) : les titres ne doivent pas être confondus avec les descriptions. Les titres sont des éléments d'identifications et sont court. Il ne devrait pas y avoir de références dans les listes.
b. Au besoin, la source doit être présentées dans la description. Mais pas dans la liste dans les pages préliminaires. Faire suivre les sources directement dans la description de l'élément
% Max : probablement le shorttitle qui est désiré   
%Method 1: Title and Description in a Single Caption (Recommended)
%\caption[Main Short Title]{\textbf{Main Title of the Figure.} \\ \small This is the secondary title or a detailed description of the graph.}

% Luc : essayé de conserver le tire sur une ligne et est un vraiment un identificateur.

% Maxime : Est-ce que le numéro de page est une limitation sur le rendu citepage dans latex, demandé, sinon mettre dans la biblio.

5. Abréviations : Voir guide

a. La définition des abréviations dans le texte doit être reprise au besoin.
% Luc : Probablement lorsque c'est réutilisé entre les chapitres au début des prochaines
b. Les abréviations de la liste devraient être uniquement celles qui ne sont pas dans les articles.
% Maxime : Besoin de plus de précisions, on prend pour acquis que le lecteur à une littérature sur les abréviations. Ceux très populaire?
% Maryse `Réponse hypothétique à la question de Maxime ... est-ce que Marie-Flavie considère les abréviations mentionnées dans les articles comme étant connues (populaires?

c. Réviser la nécessité d’introduire certaines abréviations. Il y a plusieurs abréviations qui ne sont pas nécessaires. Ça alourdit le texte pour pas grand gain de place.
% Maxime : Demander le critère pour inclure ou pas.
% Maryse : Selon les informations recueillies sur le site de l'OQLF, les abréviations apparaissant dans la liste doivent : (1) apparaitre plus d'une fois (2) être des abréviations peu connues. Ma question : qu'est-ce qui distingue une abréviation connue d'une peu connue? Son utilisation dans un article? 
d. Les abréviations doivent être en français (sauf certains acronymes ou technologies)

e. Réviser les abréviations (acronymes et technologies) en anglais et leur présentation dans le texte et dans la liste
% Luc : demander si je suis la vitrine à MF parce que je ne comprends pas exactement les instructions demandés.
% Maryse : selon ma compréhension du guide de rédaction, voici un exemple de présentation d'abréviation pouvant figurer dans la liste (cela dépend si elle est considérée connue ou pas) : CMMI  Modèle intégré de maturité des capacités (\textit{Capability Maturity Model Integration})

6. Divers

%a. Introduction : il manque le chapitre de la conclusion dans la structure du document à la fin. (fait)
b. Attention, plusieurs diagrammes sont trop petits pour que le texte soit lisible
% Luc : Mettre le titre avant et indiquer voir page suivante. MCD. Augmenter probablement la police. Changer la macro pour l'affichage horizontale.
c. Beaucoup de perte d'espace à éliminer (changements de page prématurés)


Questions et notes de Maxime :

%1. La distance entre les paragraphes (p.2) est trop petite trouver la bonne mesure. Le même problème devrait arriver au chapitre 2. (fait)
% Finalement c'était l'espace qui était trop grand.

2. Déplacer les nos de p. dans la notice (p.6), est-ce que cela signifie d'augmenter le nombre de références? (non)

3. Mettre les références dans les tableaux (p.8), cela implique forcément qu'on ne peut plus en ajouter ou en retirer parce que les numéros peuvent changer n'est-ce pas?
% MARYSE : J'ai la même compréhension. Mais à ce stade-ci, tu ne dois pas ajouter ni supprimer de références, non?

%4. La première ligne de chaque paragraphe chapitre 1 (p.13), je crois qu'il est demandé un changement de police. (fait)
% MARYSE : C'est aussi mon avis. Mais elle veut surtout que l'italique ne soit pas utilisé (Commentaire de Marie-Flavie : enlever les italiques un peu partout... (Section Langue point c.)

5. Dans (LOC, Ligne of Code en anglais), qu'est-ce qui n'est pas nécessaire?
%Maryse :  Il y a le "g" (désolée, je ne l'avais pas vu). Mais mon interprétation est que Marie-Flavie ne veut pas qu'on précise le terme anglais entre parenthèses (Section Langue, point b)

%6. Au lieu d'avoir les titres en Français entre guillemets, un changement de Police est demandé (p.76). (fait)
%Maryse : OK. (1) Prendre note que l'utilisation que j'ai faite dans ce contexte n'est pas erronée selon la Vitrine linguistique - cela ne respecte pas les normes départementales. (2) J'ai utilisé les guillemets pour mettre en évidence les termes. Il faudra utiliser une police différente pour ces termes. 
% Luc : Utiliser la police dans la figure.

%7.Le MCD, le texte trop petit. (commentaire de MF : si possible, séparer en plusieurs sous-graphes) (p.82) (p.91) (p.108), mais qu'en est-il pour p.92?  Trouver la mesure pour le texte acceptable. (fait) changé police et taille 
%Maryse : selon moi, il est aussi trop petit. Je crois que Marie-Flavie avait cessé de les annoter. La taille de la police utilisée à la figure 3.5 semble acceptable. La page 37 du document des normes présente des précisions quant aux polices et aux tailles de caractères pouvant être utilisées dans le mémoire. 
%Paratype Sans Narrow [OTF or TTF available]

8. Dans le glossaire (p.166), les parenthèses avec les termes en anglais$ sont à retirer, est-ce que c'est à faire dans l'ensemble du mémoire ou seulement dans le glossaire?
% Maryse : Dans l'ensemble du mémoire. Voir son commentaire section Langue, point b.

9. Est-ce qu'il est effectivement demandé d'inverser le contenu entre l'anglais et le français dans une liste? (p.203)
% Maryse : on dirait, mais je ne comprends pas. Marie-Flavie indique que l'utilisation de l'anglais n'est pas justifiée dans certains cas et nous invite à consulter la vitrine linguistique de l'OQLF. Mais dans le cas de citations dans une langue étrangère, nous respectons la directive : https://vitrinelinguistique.oqlf.gouv.qc.ca/23207/la-redaction-et-la-communication/bibliographie-et-citations/citations/citation-en-langue-etrangere
% Est-ce qu'il y a moyen d'obtenir plus d'infos de Marie-Flavie à ce propos?

10. Source : et etc est à placer après le titre de la figure ou du tableau et non en bas à gauche.

%11. Le titre de la section  a été renommé pour prendre un texte en début de paragraphe (p.5). Est-ce que je devrais faire partout pareil? Ou c'est seulement contextuel à ce morceau de texte comparativement aux autres sous-sections objectifs? (fait)
% Maryse : Oui, tu devrais faire cela aux sections où cela est applicable.
% Luc : le paragraphe ests réécrit.

12. Je crois qu'il est souhaité que les mots soient traduits simplement par exemple, Tableau 1.2 (p.13) éliminant les nuances linguistiques qui peut exister. Est-ce que j'ajoute une note de bas de tableau pour expliquer pourquoi cela n'est pas traduit?
% Maryse : Mes observations : (1) Ce tableau est lié à du contenu figurant à la page 12 (par.07) (2) Pas de référence bibliographique près du tableau (3) Marie-Flavie a indiqué qu'il faut utiliser davantage la terminologie francophone. Ma réponse : il faut indiquer plus clairement la source de ces termes et fournir une traduction.  

13. Le 2 c) indique je dois retirer l'italique pour ce qui est traduit. Il ne restera pas beaucoup d'italique.
% Maryse : j'ai relu le commentaire de Marie-Flavie et ce qui est mentionné dans le guide des normes. Elle mentionne "les termes en langue étrangère qui ne sont pas traduits". Selon ma compréhension, cela exclut les citations en langue étrangère (voir guide p. 38) et "et al.". Selon mon interprétation, ce commentaire vise directement les titres en gras-italique.

%14. Quelles polices devrais-je utiliser pour les titres en français et les premières lignes de paragraphe dans le chapitre 1?
% Maxime poser la question si les polices énuméré dans le guide sont pour le corps du texte.
% Maryse : voir choix possibles à la page 47 du guide des normes

15. Je crois que dans le glossaire MF m'a donnée les acronymes que je dois franciser, cela tourne pas mal autour de IA.
% : Maryse : oui, elle veut que ton mémoire soit le plus francisé possible. Cela touche la liste, le glossaire et ton mémoire.

%16. Qu'est-ce qui manque à la notice 4 ACM, ... (p.116)? Est-ce que je dois mettre le nom au complet comme pour [16], [18]?
%Maryse : Voir page 20 du guide des normes pour la liste des infos à fournir dans le cas d'une page Web.  Malheureusement, il manque l'exemple. Et ACM doit être écrit au long.
%Maxime: Poser la question puisque ISO est rarement au long.


%17. Je pense qu'il y a une confusion pour certaines notice par exemple ACM [4], ce n'est pas le document, comment je devrais l'indiquer dans la notice comme ceci? Howpublished = {Une page web}
% Maryse : voir mon commentaire à la question précédente. Plus largement, il faudrait s'assurer que tu possèdes toutes les informations nécessaires à chaque type de références, et ce, dans la bonne langue, avant de regénérer ta bibliographie.

%18. Poser la question à MF si le font et le style du gabarit sont corrects.